% !TeX root = ../main.tex
\chapter{Programmes LabVIEW (.vi)}
\label{app:labview}

Les quatre programmes LabVIEW créés pour ce laboratoire se trouvent dans le dossier \texttt{Resultat/VI/} :

\begin{itemize}[leftmargin=1.2cm]
    \item \texttt{Labo03\_NTC.vi} : mesure et conversion de température de la thermistance CTN en mode standard (cadence de \qty{20}{\milli\second}, formule inverse et interpolation par table) ;
    \item \texttt{Labo03\_NTC\_OverSampling.vi} : mesure CTN avec suréchantillonnage (moyenne de paquets de 100 échantillons) ;
    \item \texttt{Labo03\_PT100.vi} : mesure de résistance de la sonde Pt100 avec le multimètre Keysight 34460A (mode 2 fils et 4 fils, temporisation de \qty{500}{\milli\second}) ;
    \item \texttt{Labo03\_PT100\_\ AprèsEtalonage.vi} : conversion de la résistance en température avec le sous-VI RTD (paramètres $R_0$, $A$ et $B$ étalonnés) et enregistrement des mesures du saut thermique dans un fichier.
\end{itemize}
